
\documentclass{article}
\usepackage{colortbl}
\usepackage{makecell}
\usepackage{multirow}
\usepackage{supertabular}

\begin{document}

\newcounter{utterance}

\centering \large Interaction Transcript for game `textmapworld', experiment `medium\_cycle', episode 33 with reka{-}flash{-}3.1{-}t0.0.
\vspace{24pt}

{ \footnotesize  \setcounter{utterance}{1}
\setlength{\tabcolsep}{0pt}
\begin{supertabular}{c@{$\;$}|p{.15\linewidth}@{}p{.15\linewidth}p{.15\linewidth}p{.15\linewidth}p{.15\linewidth}p{.15\linewidth}}
    \# & $\;$A & \multicolumn{4}{c}{Game Master} & $\;\:$B\\
    \hline

    \theutterance \stepcounter{utterance}  
    & & \multicolumn{4}{p{0.6\linewidth}}{
        \cellcolor[rgb]{0.9,0.9,0.9}{
            \makecell[{{p{\linewidth}}}]{
                \texttt{\tiny{[P1$\langle$GM]}}
                \texttt{Please help me with the following task. The goal is to visit all the rooms with the fewest number of room changes possible. In each room, you need to decide the direction to go in. Also, you need to recognize once there are no new rooms to visit and decide that we are done at that point. Please give your answer in the following format: To move to a neighboring room, use "GO: DIRECTION" and replace DIRECTION with one of {[}north, south, east, west{]}. To stop the exploration, answer with "DONE" instead. Omit any other text.} \\
\texttt{Here is an example:} \\
\texttt{You are in the Kitchen. Currently available directions: south, west. What is your next instruction?} \\
\texttt{GO: west} \\
\texttt{You have made a step and entered a Lobby. Currently available directions: east, north. What is your next instruction?} \\
\texttt{GO: north} \\
\texttt{...} \\
\texttt{You have made a step and entered a Bedroom. Currently available directions: south. What is your next instruction?} \\
\texttt{DONE} \\
\texttt{Let us start. You are in the Billiard room. Currently available directions: north, south. What is your next instruction?} \\
            }
        }
    }
    & \\ \\

    \theutterance \stepcounter{utterance}  
    & \multicolumn{4}{p{0.6\linewidth}}{
        \cellcolor[rgb]{0.9,0.9,0.9}{
            \makecell[{{p{\linewidth}}}]{
                \texttt{\tiny{[P1$\rangle$GM]}}
                \texttt{<reasoning>} \\
\texttt{Let's tackle this problem. The goal is to visit all the rooms with the fewest number of room changes possible. So, I need to figure out the optimal path that covers all rooms in the least steps. But how do I know the layout of the rooms? The user hasn't provided a map, so maybe I need to explore step by step, keeping track of visited rooms and directions.} \\
\\ 
\texttt{Wait, the example given starts in the Kitchen, then goes west to Lobby, north to Bedroom, then done. So each move is a step, and the directions are given based on the current room. The challenge is to navigate through the rooms without getting stuck and ensuring all are visited with minimal steps.} \\
\\ 
\texttt{First, I need to keep track of which rooms I've already visited. Each time I enter a new room, I add it to a list. Then, when choosing directions, I should prioritize unvisited rooms. But if all adjacent rooms are visited, then I have to backtrack. However, backtracking might increase the number of steps, but since the goal is to visit all rooms, maybe it's necessary.} \\
\\ 
\texttt{Wait, but the problem says "the fewest number of room changes possible". Room changes are each time you move from one room to another. So each "GO" command is a room change. So the objective is to minimize the total number of GO commands needed to visit all rooms. That's essentially finding the shortest possible path that covers all rooms, which is similar to the Traveling Salesman Problem (TSP). But TSP is NP{-}hard, so for a small number of rooms, maybe a backtracking approach could work, but since the user is expecting a step{-}by{-}step solution here, perhaps a depth{-}first search (DFS) or breadth{-}first search (BFS) approach would be better.} \\
\\ 
\texttt{But since the user is simulating the exploration step by step, I need to make decisions at each step. So maybe a greedy approach: always go to the nearest unvisited room. Alternatively, use BFS to explore all possible paths and find the shortest one. But without knowing the entire map, it's impossible to precompute. So the agent has to explore dynamically.} \\
\\ 
\texttt{Wait, but the problem is presented as a text{-}based interaction where each step reveals the current room and available directions. So the assistant has to make decisions based on the current state provided by the user. Wait, no—the user is providing the current room and directions each time. Wait} \\
            }
        }
    }
    & & \\ \\

    \theutterance \stepcounter{utterance}  
    & & & \multicolumn{2}{p{0.3\linewidth}}{
        \cellcolor[rgb]{0.9,0.9,0.9}{
            \makecell[{{p{\linewidth}}}]{
                \texttt{\tiny{[GM$|$GM]}}
                \texttt{abort game} \\
            }
        }
    }
    & & \\ \\

\end{supertabular}
}

\end{document}
